% ===== German content body =====
\begin{titlepage}\centering
\vspace*{4mm}
{\displayfont\bfseries\fontsize{40}{44}\selectfont Balkon\,Borg}\\[3mm]
{\displayfont\LARGE Aufbauanleitung — Hardware}\\[1mm]
{\large\color{soft} von den Bauteilen bis zum fertigen Gerät}\\[9mm]
\dekocover{cover.png}{Titelbild: das fertige Gerät unter dem Balkon}\\[7mm]
\includegraphics[width=0.6\textwidth]{enclosure.png}\\[8mm]
\begin{tcolorbox}[width=0.86\textwidth,colback=cfun,colframe=ink,arc=4pt]\centering
Die Anleitung geht der Reihe nach vor: löten, verkabeln, einbauen, montieren. Nichts
davon ist schwierig, wenn man es Schritt für Schritt macht. Lass dir Zeit und arbeite die
Kapitel in ihrer Reihenfolge ab.
\end{tcolorbox}
\vfill
{\footnotesize\color{soft}Technische Bilder aus dem CAD/PCB des \texttt{balkon-borg}-Repos
(nicht gemalt) · Stand \today}
\end{titlepage}

\section*{Bevor du anfängst}
Dieses Heft führt durch den kompletten \textbf{Hardware-Aufbau}: Platine bestücken, Kabel
konfektionieren, alles ins Gehäuse, verdrahten, an die Decke. Die Software (ESP flashen,
Pi einrichten) ist bewusst nicht enthalten. Plane für den Aufbau in Ruhe einen ruhigen
Nachmittag ein — mit einer Pause zwischen Löten und Verdrahten.

\begin{achtung}
Das \textbf{230-V-Netzteil bleibt außen} in seinem eigenen Gehäuse. In das gedruckte
Gehäuse kommen \textbf{nur 5\,V}. Am 230-V-Teil nur arbeiten, wenn es vom Netz getrennt ist.
\end{achtung}
\begin{warum}
Kunststoff (PA12) und Netzspannung passen nicht zusammen: kein definierter Brandschutz,
keine Kriechstrecken. Deshalb sitzt das Netzteil separat, und das gedruckte Gehäuse führt
ausschließlich Kleinspannung. Das Gehäuse kommt in \textbf{zwei Hälften} aus dem Drucker
(Naht in der Mitte); beim Aufbau richten zwei 4-mm-Passstifte die Hälften aus und drei
interne M3-Klemmen verschrauben sie.
\end{warum}

\tableofcontents

% ---------------- Teileliste ----------------
\clearpage
\section{Teileliste}
\dekopikto{p-tools.png}{Teile}
\begin{ziel}
Vor dem Start alles vollständig da haben und vorsortieren. Die Zahlen sind Stückzahlen für
\emph{ein} Gerät. Verbindliche Referenz: \texttt{pcb/docs/board-spec.md} (Platine) und
\texttt{docs/cost-estimate.md} (Gesamtgerät).
\end{ziel}

\textbf{Auf der Platine (bestücken/löten):}\par
\rowcolors{2}{rowb}{white}
\begin{tabular}{clp{88mm}}
\toprule
\textbf{\#} & \textbf{Teil} & \textbf{Beschreibung} \\ \midrule
1 & ESP32-DevKitC-V4 & offizielles Espressif-Board, WROOM-32E, 38-pin, \textbf{25,4\,mm} Reihenabstand (kein 0,9"-Klon) \\
2 & Buchsenleiste 1×19 & 2,54\,mm, für den ESP (gesockelt) \\
4 & Transistor BC337 & NPN, TO-92, Treiber für die Taster-LEDs \\
2 & Widerstand 220\,$\Omega$ & Serienwiderstand Radar-UART \\
7 & Widerstand 100\,$\Omega$ & Serienwiderstand Encoder + Taster \\
4 & Widerstand 1\,k$\Omega$ & Basiswiderstand der Transistoren \\
2 & Widerstand 4,7\,k$\Omega$ & I²C-Pull-ups (optional / DNP, oft am BME-Breakout) \\
1 & Elko 10\,\textmu F & 3V3-Entkopplung (\textbf{gepolt}) \\
1 & Kondensator 100\,nF & 3V3-Entkopplung (ungepolt) \\
1 & Polyfuse ~2\,A & rücksetzbare 5-V-Sicherung am Eingang \\
1 & JST-XH Stiftleiste 2p & \texttt{J\_PWR} (5-V-Eingang) \\
7 & JST-XH Stiftleiste 4p & \texttt{J\_RADAR}, \texttt{J\_BME}, \texttt{J\_ENC}, \texttt{J\_BTN1--4} \\
\bottomrule
\end{tabular}

\textbf{Externe Baugruppen:}\par
\rowcolors{2}{rowb}{white}
\begin{tabular}{clp{88mm}}
\toprule
\textbf{\#} & \textbf{Baugruppe} & \textbf{Beschreibung} \\ \midrule
1 & Raspberry Pi 5 (8\,GB) & + Active Cooler + microSD — die zentrale Recheneinheit \\
1 & LD2410B & 24-GHz-Präsenzradar (UART), erkennt auch stillsitzende Personen \\
1 & BME280 & echter Bosch-Sensor: Temperatur / Luftfeuchte / Luftdruck (I²C) \\
1 & EC11-Encoder & Dreh-Encoder mit Taster (Helligkeit) \\
4 & Taster 12\,mm & beleuchtet, Edelstahl, \textbf{5-V}-Ring-LED \\
1 & WLED-Controller & Athom High-Power (vorgeflasht) \\
1 & SK6812 RGBW-WW & Streifen 60 LED/m: 8 Reihen à 25 = 200 px, auf Alu-Platte (438×88×3, Zuschnitt) + Opal-Diffusor \\
1 & RTL-SDR V3 & + Teleskop-SMA-Antenne (ADS-B, optional LoRa) \\
1 & USB-Mikrofon & am Pi (BirdNET) \\
1 & Audio-Kette & USB-Soundkarte (DELOCK 61645) + PAM8403-Verstärker + Visaton BF\,45 \\
1 & Kamera Modul 3 & am CSI-Port des Pi \\
1 & Mean Well LRS-150-5 & 5\,V/22\,A Netzteil, \textbf{extern} in eigenem Gehäuse \\
\bottomrule
\end{tabular}

\textbf{Kleinteile / Mechanik:}\par
\rowcolors{2}{rowb}{white}
\begin{tabular}{clp{88mm}}
\toprule
\textbf{\#} & \textbf{Teil} & \textbf{Beschreibung} \\ \midrule
~8 & Gewindeeinsatz M2.5 & Heat-Set (Pi, Trägerplatine, WLED-Halterung) \\
3 & Gewindeeinsatz M3 & Heat-Set, Nahtklemmen der beiden Gehäusehälften \\
2 & Passstift 4\,mm & richtet die Hälften an der Naht aus \\
-- & Schrauben M2.5 / M3 / M5 & M2.5 innen, M3 ×3 Nahtklemmen, M5 ×4 Deckenohren \\
1 & XT60-Buchse & 5-V-Einspeisung in der Rückwand \\
2 & Wago 221 (5-fach) & 5\,V+ und GND Sternverteilung \\
-- & Sicherungen + Halter & inline 10\,A (LED), 5\,A (Pi) \\
-- & Kabel + Aderendhülsen & 2,5 / 1,5 / 0,5\,mm² \\
-- & JST-XH Crimp-Set & Kontakte + Gehäuse für die Sensor-/Bedienkabel \\
-- & Diverses & Panel-Muttern, Thermal-Pad, Kleber, Insektengitter, Kabelbinder \\
\bottomrule
\end{tabular}

\begin{schmunzel}
Sortier die Widerstände gleich nach Wert in kleine Schälchen. Vier verschiedene Werte
später am Bauteilgurt zu unterscheiden, kostet mehr Nerven als der ganze Rest.
\end{schmunzel}

% ---------------- Werkzeug ----------------
\clearpage
\section{Werkzeug}
\dekopikto{p-solder.png}{Werkzeug}
\begin{itemize}[leftmargin=8mm,itemsep=2pt]
  \item \faBolt\ temperaturgeregelter Lötkolben (feine Spitze) + Lötzinn, Entlötlitze
  \item \faCut\ Seitenschneider, Abisolierzange, \textbf{JST-XH-Crimpzange}
  \item \faScrewdriver\ kleiner Kreuzschlitz-Schraubendreher, Pinzette
  \item \faFire\ Lötkolben oder Heißluft für die \textbf{Gewindeeinsätze}
  \item \faSearch\ Multimeter (Durchgangsprüfer) für den Sicht-Check am Ende
\end{itemize}
\begin{tipp}
Ein 30-Euro-Lötkolben mit Temperaturregelung reicht völlig (~330\,°C für bleifreies Zinn).
Die JST-Crimpzange ist die einzige Spezialzange und spart am meisten Zeit.
\end{tipp}

% ---------------- A Platine ----------------
\clearpage
\section{Platine bestücken (löten)}
\dekopikto{p-solder.png}{löten}
\begin{ziel}
Aus der leeren Platine wird die fertig bestückte Trägerplatine mit aufgestecktem ESP32.
Bestückt wird \textbf{von flach nach hoch}, damit beim Umdrehen immer alles plan aufliegt.
\end{ziel}
\begin{schritte}
  \item \textbf{15 Widerstände} (R1–R15) einsetzen: alle stecken, Beinchen leicht umbiegen,
        dann in einem Durchgang verlöten und die Enden kürzen. Widerstände sind ungepolt —
        egal herum.
  \item \textbf{4 Transistoren} BC337 (Q1–Q4) einsetzen — \emph{auf die Orientierung achten}
        (siehe Achtung-Box).
  \item \textbf{Kondensatoren}: C1 (10\,\textmu F Elko, \textbf{gepolt} — der Minus-Streifen
        gehört auf die GND-Seite) und C2 (100\,nF, ungepolt).
  \item \textbf{Polyfuse F1} einlöten (die 5-V-Sicherung am Eingang).
  \item \textbf{8 JST-XH-Stiftleisten} einlöten: \texttt{J\_PWR} (2p) sowie \texttt{J\_RADAR},
        \texttt{J\_BME}, \texttt{J\_ENC}, \texttt{J\_BTN1--4} (je 4p).
  \item \textbf{2 Buchsenleisten} (1×19) für den ESP — vor dem Verlöten gerade ausrichten.
  \item \textbf{ESP32-DevKitC-V4 aufstecken}.
\end{schritte}
\bild{pcb-top.png}{0.82\textwidth}{Die bestückte Platine mit Bauteil-Bezeichnern (Quelle: KiCad-Render).}
\begin{achtung}
BC337: die \textbf{flache Seite} des Transistors muss zur Markierung auf der Platine
(Silk) zeigen. Falsch herum schalten die Taster-LEDs nie. Erst schauen, dann löten. Und:
den Elko C1 richtig herum — der Minus-Streifen zeigt weg von der 3V3-Markierung.
\end{achtung}
\begin{tipp}
Buchsenleisten kippen gern schief. Steck den ESP \emph{vor} dem Verlöten kurz auf, richte
alles rechtwinklig aus, löte dann zuerst nur je einen Pin außen und korrigiere, bevor du
den Rest lötest.
\end{tipp}
\begin{ergebnis}
Eine vollständig bestückte, DRC-saubere Trägerplatine mit sitzendem ESP32.
\end{ergebnis}

% ---------------- B Kabel ----------------
\clearpage
\section{Kabel konfektionieren (crimpen)}
\dekopikto{p-crimp.png}{crimpen}
\begin{ziel}
Alle Verbindungskabel von der Platine zu Sensoren, Bedienteilen und Stromversorgung — mit
der richtigen Pinbelegung und Aderfarbe. \textbf{Wichtig:} die JST-Kabel werden
\emph{gecrimpt}, nicht gelötet. Gelötet wird nur, wo ein Bauteil offene Pads statt eines
Steckers hat.
\end{ziel}
\begin{schritte}
  \item Kabellängen grob abmessen (großzügig, im Gehäuse ist Platz).
  \item Adern ~3\,mm abisolieren und die \textbf{JST-XH-Kontakte crimpen} (Crimpzange).
  \item Kontakte ins Steckergehäuse schieben, bis sie hörbar einrasten.
  \item Aderfarbe konsequent: \ader{wRD}{5\,V} \quad \ader{wBK}{GND} \quad
        \ader{wOG}{3,3\,V} \quad \ader{wGN}{Signal} \quad \ader{wYE}{Signal} \quad
        \ader{wBU}{Signal} \quad \ader{wWH}{Signal}.
\end{schritte}
\bild{wiring-harness.png}{0.78\textwidth}{Der komplette Kabelbaum: welches Kabel an welchen Pin, mit Aderfarben (Quelle: Netzliste → WireViz).}
\begin{warum}
Konsequente Aderfarben (rot = 5\,V, schwarz = GND) sind kein Selbstzweck: sie verhindern,
dass du beim Zusammenstecken Plus und Minus vertauschst — der häufigste und ärgerlichste
Fehler, mit Farbcode kaum möglich.
\end{warum}

% ---------------- C Sensoren + Verbindungstabellen ----------------
\clearpage
\section{Sensoren \& Bedienteile anschließen}
\dekopikto{p-connect.png}{anstecken}
\begin{ziel}
Jedes Kabel bekommt sein Bauteil — Ader für Ader. Unten steht für jeden Stecker genau,
welcher Pin über welche Aderfarbe an welchen Bauteil-Pin geht.
\end{ziel}

\bild{board-annotated.png}{0.86\textwidth}{Welcher Stecker zu welchem Bauteil führt (Board-Render mit Callouts).}

\textbf{Beispiel als Schema — Radar (\texttt{J\_RADAR} → LD2410B):}
\begin{center}\begin{tikzpicture}
  \node[wpin,fill=ink,text=white] at (0,0.6) {J\_RADAR};
  \node[wpin,fill=ink,text=white] at (7,0.6) {LD2410B};
  \wire{1}{1 · 5V}{wRD}{VCC}
  \wire{2}{2 · GND}{wBK}{GND}
  \wire{3}{3 · RX-Sig.}{wGN}{RX}
  \wire{4}{4 · TX-Sig.}{wYE}{TX}
\end{tikzpicture}\end{center}

\textbf{Alle Steckerbelegungen (Board-Pin — Ader — Bauteil-Pin):}\par
\rowcolors{2}{rowb}{white}
\begin{tabular}{llll}
\toprule
\textbf{Stecker} & \textbf{Board-Pin} & \textbf{Ader} & \textbf{Bauteil-Pin} \\ \midrule
\texttt{J\_RADAR} & 1 · 5V & \ader{wRD}{} & LD2410B VCC \\
 & 2 · GND & \ader{wBK}{} & GND \\
 & 3 · RX-Signal & \ader{wGN}{} & RX \\
 & 4 · TX-Signal & \ader{wYE}{} & TX \\
\texttt{J\_BME} & 1 · 3V3 & \ader{wOG}{} & BME280 VCC \\
 & 2 · GND & \ader{wBK}{} & GND \\
 & 3 · SCL & \ader{wWH}{} & SCL \\
 & 4 · SDA & \ader{wGN}{} & SDA \\
\texttt{J\_ENC} & 1 · A & \ader{wGN}{} & EC11 A \\
 & 2 · B & \ader{wYE}{} & B \\
 & 3 · SW & \ader{wBU}{} & Taster \\
 & 4 · GND & \ader{wBK}{} & GND \\
\texttt{J\_BTN1--4} & 1 · SW & \ader{wGN}{} & Taster-Schalter \\
 & 2 · GND & \ader{wBK}{} & Schalter-COM \\
 & 3 · 5V & \ader{wRD}{} & LED\,+ (Anode) \\
 & 4 · LED-K & \ader{wWH}{} & LED\,-- (Kathode) \\
\texttt{J\_PWR} & 1 · 5V & \ader{wRD}{} & Netzteil 5\,V (Wago) \\
 & 2 · GND & \ader{wBK}{} & Netzteil GND (Wago) \\
\bottomrule
\end{tabular}
\begin{tipp}
Kauf für die Taster die \textbf{5-V-Variante} der Ring-LED — dann braucht der Treiber auf
der Platine keinen Extra-Widerstand. Hat ein Bauteil offene Lötpads statt eines Steckers
(manche Radar-/BME-Breakouts), löte die Ader dort direkt an — Reihenfolge und Farbe bleiben
gleich wie in der Tabelle.
\end{tipp}
\begin{ergebnis}
Alle vier Sensor-/Bedien-Baugruppen hängen korrekt belegt an ihren Kabeln.
\end{ergebnis}

% ---------------- D Inserts ----------------
\clearpage
\section{Gehäuse vorbereiten: Einsätze setzen, Hälften fügen}
\dekopikto{p-insert.png}{Insert setzen}
\begin{ziel}
Die Gewindeeinsätze sitzen, die beiden Gehäusehälften sind verschraubt — und das LED-Panel
liegt schon im Gehäuse, denn nach dem Fügen passt es durch keine Öffnung mehr.
\end{ziel}
\begin{schritte}
  \item \textbf{M2.5-Einsätze} für Pi 5, Trägerplatine und WLED-Halterung setzen (~8 Stück).
  \item \textbf{M3-Einsätze} in die drei Nahtklemmen einer Hälfte (2× Rückwand, 1× vorn
        im Boden). Die \textbf{Deckenohren} haben M5-Durchgangslöcher — kein Einsatz nötig.
  \item \textbf{4-mm-Passstifte} in die Dübelposten einer Hälfte (ein Tropfen Kleber hält sie).
  \item \textbf{LED-Panel einlegen, dann fügen:} das Panel samt Aluplatte beim
        Zusammenschieben der Hälften mittig einlegen und dabei etwa 2\,cm hinter der
        Frontkante halten (dort kollidiert es weder mit der vorderen Nahtklemme noch mit
        den Positioniernasen). Hälften auf die Stifte schieben und mit den drei
        \textbf{M3-Schrauben} verschrauben — von der offenen Front aus, langer
        Schraubendreher.
  \item Das Panel danach \textbf{nach vorn} zwischen die vier Positioniernasen schieben
        (verklebt wird es erst später, mit dem Diffusor).
\end{schritte}
\bild{half-left.png}{0.58\textwidth}{Ansicht einer Gehäusehälfte (Innenraum sichtbar).}
\begin{achtung}
Einsätze mit dem Lötkolben \textbf{gerade und nicht zu heiß} einsenken. Schief eingesetzt
reißt später das Gewinde aus — lieber langsam und senkrecht.
\end{achtung}
\begin{ergebnis}
Alle Montagepunkte haben feste M2.5-Gewinde.
\end{ergebnis}

% ---------------- E Einbau ----------------
\clearpage
\section{Komponenten einbauen}
\dekopikto{p-install.png}{einbauen}
\begin{ziel}
Alle Baugruppen sitzen an ihrem Platz — auf ~8\,mm Abstandshaltern, mit reichlich Luft zum
LED-Panel.
\end{ziel}
\begin{schritte}
  \item \textbf{Raspberry Pi 5} an die Rückwand-Bosse schrauben.
  \item \textbf{Trägerplatine} in die rechte Bucht auf ihre M2.5-Bosse (Bauteile zur Front).
  \item \textbf{LED-Panel} (liegt seit dem Fügen im Gehäuse, auf seiner Aluplatte als
        Kühlkörper) vorn an den Rahmenrücken zwischen die Nasen kleben. Den
        \textbf{Opal-Diffusor} erst ganz zum Schluss in den Falz kleben — danach ist der
        Innenraum nur noch über die Deckenöffnung erreichbar.
  \item \textbf{Turm-LEDs + Radar} in den nach unten ragenden Turm; das Radar von innen
        über den 30×5,5-Schlitz kleben (Kleberaupe rundum — hält und dichtet gegen Insekten).
  \item \textbf{Kamera} in die hängende Box: Objektivblock flach an die geneigte Frontwand,
        Tubus durch die Bohrung, Platinenecken mit \textbf{Heißkleber} fixieren. Dazu
        \textbf{Mikrofon} und \textbf{Audio} (USB-Soundkarte, PAM8403-Verstärker,
        Visaton-Lautsprecher).
  \item \textbf{RTL-SDR} an einen \emph{USB-2}-Port, Antenne an die SMA-Buchse.
\end{schritte}
\begin{figure}[htbp]\centering
  \includegraphics[width=0.47\textwidth]{view-inside.png}\hfill
  \includegraphics[width=0.47\textwidth]{view-bottom.png}\par
  {\footnotesize\color{soft}Innenraum (links) und Unterseite mit Kamera-/Turmöffnung (rechts).}
\end{figure}
\begin{warum}
Der SDR gehört an USB-2, nicht USB-3: USB-3 strahlt genau auf 1090\,MHz und würde den
Flugzeug-Empfang (ADS-B) stören. Kleiner Handgriff, großer Unterschied.
\end{warum}
\begin{ergebnis}
Ein voll bestücktes Gehäuse — nur noch nicht verdrahtet.
\end{ergebnis}

% ---------------- F Verdrahtung ----------------
\clearpage
\section{Stromversorgung \& Verdrahtung}
\dekopikto{p-wire.png}{verdrahten}
\begin{ziel}
Die 5-V-Sternverteilung steht (mit Sicherungen), und alle Kabel von der Platine zu den
Sensoren sind gesteckt.
\end{ziel}
\begin{schritte}
  \item \textbf{XT60-Buchse} in die Rückwand, 5\,V vom externen Netzteil hierher.
  \item \textbf{Wago-Verteiler} für 5\,V+ und GND (gemeinsame Masse für alles).
  \item \textbf{Sicherungen} inline: 10\,A (LED-Panel), 5\,A (Pi\,5), 2\,A (Platine über deren F1).
  \item Zweige mit passendem Querschnitt: 2,5 / 1,5 / 0,5\,mm².
  \item \textbf{Audio-Verstärker} 5\,V/GND vom \emph{Pi-Zweig} abgreifen (gleiche Masse → kein Brummen).
  \item Alle \textbf{JST-Kabel} Platine\,↔\,Sensoren stecken — Belegung nach der Tabelle in Kapitel~5.
\end{schritte}
\begin{achtung}
Vor dem ersten Einschalten: \textbf{Polung prüfen}. 5\,V verpolt anzulegen killt Pi und
Platine sofort. Rot = 5\,V, Schwarz = GND, überall.
\end{achtung}
\begin{tipp}
Halt die 5-V-Leitung kurz und dick, besonders zum LED-Panel — sonst fällt am Strangende die
Helligkeit ab. Das lange Kabel gehört auf die 230-V-Seite (zum externen Netzteil).
\end{tipp}
\begin{ergebnis}
Ein verdrahtetes Gerät mit abgesicherter Stromversorgung und gemeinsamer Masse.
\end{ergebnis}

% ---------------- G Montage ----------------
\clearpage
\section{Deckenmontage}
\dekopikto{p-assemble.png}{montieren}
\begin{ziel}
Das Gehäuse hängt sicher unter dem Balkon. Die Hälften sind seit dem Vorbereiten
verschraubt — hier geht es nur noch an die Decke.
\end{ziel}
\begin{schritte}
  \item Letzter \textbf{Kabel-Check}: nichts eingeklemmt, nichts geknickt.
  \item Das Gerät über die \textbf{Deckenohren} (M5, ×4) von unten an die Decke schrauben —
        eine Unterlegscheibe unter jedem Kopf hilft.
\end{schritte}
\begin{warum}
Die vier Ohren tragen das gesamte Gewicht. Nimm Dübel/Anker passend zum Deckenmaterial;
im Zweifel eine Nummer größer.
\end{warum}
\begin{ergebnis}
Ein geschlossenes, montiertes Balkon-Borg an seinem Platz.
\end{ergebnis}

% ---------------- H Check ----------------
\clearpage
\section{Sicht-Check vor dem Einschalten}
\dekopikto{p-check.png}{prüfen}
\begin{ziel}
Einmal alles durchgehen, bevor Strom drauf kommt. Die Software ist ein eigenes Kapitel —
hier geht es nur um die Hardware.
\end{ziel}
\begin{schritte}
  \item Polung aller 5-V-Anschlüsse geprüft (rot/schwarz)?
  \item Keine Lötbrücken/Kurzschlüsse an der Platine (Durchgangsprüfer)?
  \item Alle JST-Stecker vollständig eingerastet, Belegung wie in Kapitel~5?
  \item Sicherungen (10/5\,A) an ihrem Platz?
  \item Das 230-V-Netzteil ist außen und noch \textbf{vom Netz getrennt}?
\end{schritte}
\begin{ergebnis}
Ein aufgebautes, geprüftes Gerät — bereit für die Software (anderes Heft).
\end{ergebnis}
